\section{Performance Evaluation}
\label{sec:performance}

This section provides the performance analysis and benchmark data required by the call. The current repository uses repository-local measurements for the packaged Kylin NGCC instances.

\paragraph{Performance analysis.}
At the design level, Kylin's dominant costs are hash and XOF calls. Key generation builds the top hypertree layer, signing computes BPORS+FP material and one authentication path through each hypertree layer, and verification recomputes the BPORS+FP public key and verifies the hypertree signature. \WOTSCS reduces the number of $n$-byte elements in each one-time signature, BPORS+FP reduces the bottom few-time signature payload, and Diverse Layer Heights keeps the signing hash-call proxy balanced even when $h$ is not divisible by $d$.

\begin{table}[H]
    \centering
    \caption{Design-level hash-call and signature-size metrics for the primary Kylin instances.}
    \label{tab:kylin-hashcall-performance}
    \scriptsize
    \begin{tabular}{l r r r r}
    \toprule
    Instance & KeyGen hash calls & Sign hash calls & Verify hash calls & Signature bytes \\ \midrule
    Kylin-128 & 417279 & 2389373 & 7362 & 6472 \\
    Kylin-256 & 324095 & 3382025 & 11405 & 25352 \\
    Kylin-512 & 519295 & 6493104 & 44418 & 91528 \\
    \bottomrule
    \end{tabular}
\end{table}

\paragraph{Benchmark platform.}
The software implementation is written in ISO C and is compiled with GCC 13.3.0. The reference implementations keep their default \texttt{-O2} ISO C build with warning checks enabled. The optimized 64-bit PC implementations use per-directory builds centered on \texttt{-O3 -march=native}; \texttt{Kylin-128} additionally enables link-time optimization, omitted frame pointers, loop unrolling, and \texttt{-fno-plt}. The benchmark host is an x86-64 64-bit PC environment: AMD Ryzen 7 7800X3D 8-Core Processor, 8 physical cores and 16 hardware threads, approximately 4.192 GHz as reported by \texttt{/proc/cpuinfo}, 30 GiB RAM, Ubuntu 24.04.3 LTS on Linux 6.18.33.1-microsoft-standard-WSL2. The implementation uses the ICCS auxiliary hash/XOF interfaces \texttt{sm3hash}, \texttt{pseudohash}, and \texttt{pseudoXOF}; no external cryptographic library is required for the reported package benchmarks.

\paragraph{Software performance, transmission, and storage cost.}
The automated package measurements in \cref{tab:kylin-package-bench} are produced from the current implementation tree by the package benchmark driver. The table reports the median of 9 measured host-local wall-clock timings on a 32-byte message after 3 warm-up iterations, and includes the public-key, secret-key, and signature sizes in bytes. Because these are repository-local wall-clock measurements, the absolute values should be interpreted as host-specific data points rather than cross-platform claims.

In the current rerun, the optimized implementations remain faster than the
reference ones, although the retuned BPORS+FP parameters change the relative
signing gains. For \texttt{Kylin-128}, \texttt{Kylin-256}, and
\texttt{Kylin-512}, the measured signing speedups are approximately
$3.39\times$, $1.30\times$, and $2.60\times$, respectively; the key-generation
speedups are approximately $3.92\times$, $4.27\times$, and $6.10\times$; and
the verification speedups are approximately $2.59\times$, $2.66\times$, and
$3.93\times$.

\begin{table}[H]
    \centering
    \caption{Repository-local software performance and transmission/storage cost for the packaged Kylin NGCC instances. Times are milliseconds on a 32-byte message; sizes are bytes.}
    \label{tab:kylin-package-bench}
    \scriptsize
    \begin{tabular}{l c c c c c c c}
    \toprule
    Instance & Impl. & KeyGen (ms) & Sign (ms) & Verify (ms) & Pk & Sk & Sig \\ \midrule
    \csvreader[head to column names]{paper/data/package_bench.csv}{}%
    { \name & \impl & \keygenms & \signms & \verifyms & \pkbytes & \skbytes & \sigbytes \\ }
    \vspace{-1em} \\ \bottomrule
    \end{tabular}
\end{table}
